\documentclass[11pt]{article}
\usepackage[a4paper,margin=22mm]{geometry}
\usepackage{fontspec}
\setmainfont{LinLibertine_R.otf}[BoldFont=LinLibertine_RB.otf,ItalicFont=LinLibertine_RI.otf,BoldItalicFont=LinLibertine_RBI.otf]
\usepackage{xcolor}
\usepackage{hyperref}
\hypersetup{colorlinks=true,urlcolor=blue}
\setlength{\parindent}{0pt}
\setlength{\parskip}{6pt}
\sloppy
\emergencystretch=2em
\begin{document}
\section*{Historical Britannica: extracted source}
Encyclopaedia Britannica - Eighth edition, Volume 6, Burning glasses-Climate - EB.16\par
{\small Source row 2292. Parser: nvidia/nemotron-3.5-parse.\par}
{\small This is model-extracted text, not a corrected transcription. OCR errors, reading order and region alignment require human review.\par}
\hrule\medskip
{\small\itshape 508\par}

{\scriptsize\color{gray} Region region-0; provider order 0.\par}
\medskip

\subsection*{CHEMISTRY.}

{\scriptsize\color{gray} Region region-1; provider order 1.\par}
\medskip

Chamistry, and so on. These basic are very much similar to the others. They are called inside bases, and the general milk is NH2 (C<sub>1</sub>H<sub>3</sub>+)). It will be seen by the first column, that several of them are known with amide bases, but the second third columns, which give the rational formula; show that they are distinct in constitution as they are in properties, although analogous. Human becomes one that the general formula of the site series (1), 2(C, H<sub>4</sub>+) may also be written (C<sub>1</sub>H<sub>4</sub>+) which is a name and that of the preceding section is its second form. The same general form also (C<sub>1</sub>H<sub>4</sub>+) N, includes also the next series, although it has may be expressed so as to show the differencee
\par

{\scriptsize\color{gray} Region region-2; provider order 2.\par}
\medskip

\textbackslash{}begin\{tabular\}\{cc\}
Ammonia \& N:N\textbackslash{}\textbackslash{}
Trankybianine \& C-S(N\_A)\textbackslash{}\textbackslash{}
Tristeklybianine \& C-S(N\_A)\textbackslash{}\textbackslash{}
Tristeklybianine \& C-S(N\_A)\textbackslash{}\textbackslash{}
Tristeklybianine \& C-S(N\_A)\textbackslash{}\textbackslash{}
Mobykalohiyianina \& C-S(N\_A)\textbackslash{}\textbackslash{}
Bipi Moharykainne \& C-S(N\_A)\textbackslash{}\textbackslash{}
Methylisht yanamine \& C-S(N\_A)\textbackslash{}\textbackslash{}
\textbackslash{}end\{tabular\}
\par

{\scriptsize\color{gray} Region region-3; provider order 3.\par}
\medskip

Even among the few here given, it will be seen that there are some kindness bases in the last table, and two of the data also with two bases in the first table, series IX. But the rational formula in the last two columns show that they differ in constituting, as they do also in properties. The base of this series are called tricylle bases.
\par

{\scriptsize\color{gray} Region region-4; provider order 4.\par}
\medskip

The three series just mentions of all of which ammo- nis is the type, include all the known \_volatile organic\_ bases, nutrient or artificial, that is, all such as are not decomposed but volatilised, unclaged by a due application of heat. An amide base is easily converted into an inside base by the acting of the iodide of one of the radicals of series I, and an inside base is converted by the measure insulin a nitryle base. J, and all of them are for delivery from anemia by the action, in successive steps, of such an iodide. Let us take iodide of ethyla, AcI. The first action is NH1+ AcI +NH1A,H1 That Is, the iodine takes 1 eq of hydrogen, forming hydroxide acid, while the ethylene replaces that as of hydrogen, forming the amino acids, synthetic which combines with theabolonic acid. The next step is the action of iodide of ethylen or ethylamin, NH1, Ae+Ac I LNH1A, HI1, the result is ivhydrate diethylaminas. And in the third stage, when iodide of ethylb acts or diethylsilylte volume in hydroxide of diethylaminas, NH Ac+AcI=N AAC, HT. Poyaling can show more mainly that if these bases are simply ammonia, by hydrogen of which is replaced in part of an whole, by positive malaria of series I. It is also evident, that since there are only 3 eps, of hydrogen in ammonia, we cannot push further the substitution of these radicals in hydrogen, as we have replaced the whole 3 equivalents. If we attempt to do so, we find that we obtain natural new bases, but of a different type, being no longer volatile, but decompound by heat, and containing oxygen. The explanation of this is very interesting. We cannot indeed replace more than three atoms of hydrogen in ammonia by our analysis, but we can aid to the 3 egs of malola a fourth, and no more, etc. in other words, we can replace the fourth eq. of hydrogen also in \_ammonium\_, which contain 4 egs of hydrogen. This leads to the formation of the next set.
\par

{\scriptsize\color{gray} Region region-5; provider order 5.\par}
\medskip

18. The tightened series is a non removable and intercepted one, of fixed lines, that is, however which cannot be volatilized without decomposition, and which no longer are of the type ammonia, but, in composition of the type of hydrated oxide of ammonium, and at the type is unknown in a separate form, we can compare them to nothing as well as to hydrated ions of potassium, or cause protein, which they are in a most astonishing degree analogue. To illustrate their information, like to suppose iodide of ethylb to act to triethylamine. Here there is no more hydrogen left to be replaced, for all 8 egs, are already replaced, and all the hydrogen present is in the form of ethylc. The action is N Ac¡+Atel=N Aac¡L. The ethylene of the iodide
\par

{\scriptsize\color{gray} Region region-6; provider order 6.\par}
\medskip

17. The seventeenth series is again one of volatile bases, Chemistry, of the type of ammonia, and of the general formula (C<sub>1</sub>H<sub>4</sub>+) N, regarded the this with the two preceding series, but has just been stated. But the more precise formula here is a function is taken from ammonia, revealed, the as amide 1 eq, hydrogen, nor as imide 4-2 eqs, hydrogen, but as nitro- gen+ 9 eqs, hydrogen, 1⁄4 Ef). This being the type, in base based the whole 3 eqs, of hydrogen are replaced by methylus or by ethyle, lec, or partly by one and partly by another of the radicals of series I. The following are some of them. Their general formula, strictly noted, is N+3 (CsI HAs+)
\par

{\scriptsize\color{gray} Region region-7; provider order 7.\par}
\medskip

unites with triethylamine or form a compound metal, tetre- yliums Aeq, analogous tamammonium NH, and thin metal have with the iodine, forming the iodide of tetrehylum, N Aa¡L is exactly similar to iodide of ammonium NH. To mention the below the iodide, which is a crystallizable salt, is acted in by oxide of silver and water. The action is, NAx¡+AtøO+Ho+Atj+NAgV,HO and the re- sults any, iodide of silver and hydrated oxide of tetrehylum. The latter is obtained as a crystalline mass by evaporation to rocorns.
\par

{\scriptsize\color{gray} Region region-8; provider order 8.\par}
\medskip

In common, it corresponds to the hydrated oxides of ammonium and potassium. If have—
\par

{\scriptsize\color{gray} Region region-9; provider order 9.\par}
\medskip

Hybrided oxide of ammonium... NH(O) HO = Am0,IO
\par

{\scriptsize\color{gray} Region region-10; provider order 10.\par}
\medskip

-- poisonim... KO IO
\par

{\scriptsize\color{gray} Region region-11; provider order 11.\par}
\medskip

(Thin NaAc, tetrehylum). Hydrated oxide of ammo- nism is not known to rise was attempting to form the radical ammonium, it is naturally resolved into ammonia and water. NH(O) HO NH+2-2RO. The reason of this would seem to be, that the attraction of the oxygen of the oxide for the fourth eq. of hydrogen is so strong as to detect is, that eq being of course and by a residue force than the others. If however, this hydrated oxide could ex- in a separate form—and possibly by some day enc- used in obtaining, perhaps under the influence of intense cold-st-would certainly resemble, caustic potash, and the hybrated oxide of tetrehylum. The latter circle in a separate form, just because the attraction of oxygen for the fourth eq. of ethylb say much less than the hydrogone that that is, as certain temperatures [for at higher temperatures by- dried oxide of tetrehylum undergoes precisely the same changes that of ammonium at certainity temperatures being converted into ethylamin, corresponding to ammonia, and hydrated oxide of ethyle, corresponding to 2 eqs of water. NaAcO,Ho+NaAc+AtO,HO. His on this account that we conjecture that under interestoid the ammonium component might be intrinsicment. The action of heat on this hybrid oxide of tetrehylum is characteristic of this whole series of bases.
\par

{\scriptsize\color{gray} Region region-12; provider order 12.\par}
\medskip

To properties, hydrated oxide of tetrehylum is so very like caustic potash, that were it not for the action of heat, in might, in addition, both confounded with it. It has an exotic alkaline taste, feels soapy to the fingers, as potash does, attracts water and carbon acid from the atmosphere, pre- epittive the salts of metals in potash does, converts oils into easy work boiled with them, and forms insoluble or sparingly soluble salts, with herbicarbonate and other tests, exactly as potash does. The chief difference is, that besides the counter alkalinate state, it has a bitter taste; and this is not of the whole class of bases of which it belongs. This and the action of heat distinguish them from potash and soda. But those some characters connect them with the fixed vegetable bases, such asquining morphins, crych-
\par

{\scriptsize\color{gray} Region region-13; provider order 13.\par}
\medskip
\newpage\section*{Source record}
bri-2292-93b4d1f3ff8c\par
Image SHA-256: \texttt{93b4d1f3ff8c91cea76506298710cd1e}\par\texttt{f022efe47dc5ba2ebe8650a7496fc589}\par
Dataset revision: \texttt{309f765b1ecc83e0a7f7}\texttt{762e44b4a927a7555443}\par
National Library of Scotland, Encyclopaedia Britannica Illustrated (CC0).\par
\url{https://huggingface.co/datasets/NationalLibraryOfScotland/encyclopaedia_britannica_illustrated/viewer/default/train?row=2292}
\end{document}
